% Originally: content/week-02.tex, \section{Completeness} (Corsin Week 2, Friday)

\section{Completeness}
\label{sec:completeness}

\begin{definition}[Cauchy sequence and complete metric space]
\label{def:complete_metric_space}
Let $(X,d)$ be a metric space and $(x_n)_{n=0}^{\infty}$ a sequence in $X$. We say that
$(x_n)_{n=0}^{\infty}$ is \newterm{Cauchy} if
\[\forall \varepsilon > 0 \ \exists N \in \mathbb{N} \text{ such that } d(x_n, x_m) < \varepsilon \quad \forall n, m \geq N.\]
$(X,d)$ is a \newterm{complete metric space} (\germanterm{vollst\"andiger metrischer Raum}) if
every Cauchy sequence in $X$ converges with respect to $d$.
\end{definition}

% Supplement: Sascha Brack/Class Notes/Week_02_Notes_Monday_Updated.pdf, pp. 7--9
\begin{proposition}[Convergence in $\mathbb{R}^n$ is coordinatewise; $\mathbb{R}^n$ is complete]
\label{prop:Rn_coordinatewise_complete}
Equip $\mathbb{R}^n$ with the standard metric $d_2$, and write
$x_m = (x_{m,1},\dots,x_{m,n})$.
\begin{enumerate}[label=\textbf{(\alph*)}]
  \item $x_m \to x$ in $\mathbb{R}^n$ if and only if $x_{m,i} \to x_i$ in $\mathbb{R}$ for every
    $i = 1,\dots,n$.
  \item $(\mathbb{R}^n, d_2)$ is complete.
\end{enumerate}
\end{proposition}

\begin{proof}
Everything follows from the two-sided estimate, valid for every $j$:
\[|x_{m,j}-x_j| \ \leq\ \lVert x_m-x\rVert \ =\ \sqrt{\sum_{i=1}^n|x_{m,i}-x_i|^2}
  \ \leq\ \sqrt{n}\,\max_i|x_{m,i}-x_i|.\]

\textbf{(a)} If $\lVert x_m-x\rVert \to 0$, the left inequality forces each coordinate to
converge. Conversely, if every coordinate converges then $\max_i|x_{m,i}-x_i| \to 0$ (a maximum
of \emph{finitely} many null sequences), so the right inequality forces
$\lVert x_m - x\rVert \to 0$.

\textbf{(b)} Let $(x_m)$ be Cauchy in $\mathbb{R}^n$. The same left inequality gives
$|x_{m,j}-x_{k,j}| \leq \lVert x_m-x_k\rVert$, so each coordinate sequence $(x_{m,j})_m$ is
Cauchy in $\mathbb{R}$. Since $\mathbb{R}$ is complete --- this is the completeness axiom from
\emph{Analysis I}, and the only input the proof has --- each converges, say
$x_{m,j} \to x_j$. Setting $x := (x_1,\dots,x_n)$, part \textbf{(a)} gives $x_m \to x$.
\end{proof}

\begin{ainote}
Worth noting how little is proved here: completeness of $\mathbb{R}^n$ is completeness of
$\mathbb{R}$, applied $n$ times. That is also where the argument would fail in infinite
dimensions --- $\max_i$ over infinitely many coordinates need not tend to $0$ even when each
does, which is exactly the phenomenon behind the $\varepsilon$-net argument in
\Cref{lem:sequentially_compact_totally_bounded}.
\end{ainote}

% Source: Corsin Nick/Class Notes/Week 2.pdf, pp. 1--13
\begin{example}[Completeness examples]
\begin{itemize}
  \item The Euclidean spaces $(\mathbb{R}^n, \langle\cdot,\cdot\rangle_{\text{eucl}})$ are
    \textbf{complete}.
  \item Any \textbf{non-closed} subset $U \subseteq \mathbb{R}^n$ (standard metric) and the
    rationals $\mathbb{Q}$ are \textbf{not complete}.
  \item $\big(C^0([0,1]), \sup_{x \in [0,1]}|\cdot|\big)$ is \textbf{complete} (Zorich, p.\ 22).
\end{itemize}
\end{example}

% Generator: Claude Opus 5 (Medium)
\begin{remark}[Completeness is not the same as closedness]
\label{rem:complete_vs_closed}
The second bullet is easy to over-read as \qt{complete means closed}. It does not, and the
difference is worth pinning down now (because it recurs).

\textbf{Closedness is relative, completeness is intrinsic.} \qt{Closed} only means anything once
you say \emph{closed in what} (i.e., it is a statement about how a set sits inside an ambient space).
Completeness asks nothing about an ambient space --- only whether the Cauchy sequences of the
space converge \emph{within it}. So:
\begin{itemize}
  \item $\mathbb{Q}$ is \textbf{closed in $\mathbb{Q}$} --- every space is closed in itself
    (\cref{ex:X_empty_clopen}) --- and yet \textbf{not complete}: the sequence
    $1,\,1.4,\,1.41,\,1.414,\dots$ is Cauchy in $\mathbb{Q}$ and converges to nothing there.
  \item $(0,1)$ is \textbf{not closed in $\mathbb{R}$} and also not complete, which is the
    second bullet's case --- but the two failures have different causes, and the first does not
    imply the second in general.
\end{itemize}

What \emph{is} true is a one-directional statement, and only inside a complete ambient space:
if $X$ is complete and $A \subseteq X$, then
\[A \text{ closed in } X \iff A \text{ complete}.\]
This is why the second bullet is correct \emph{as stated for $\mathbb{R}^n$} --- it silently uses
completeness of the ambient $\mathbb{R}^n$. Drop that hypothesis and the equivalence dies, as
$\mathbb{Q}$ inside $\mathbb{Q}$ shows.

\begin{ainote}
The same relative-versus-intrinsic distinction is the one behind
\Cref{rem:common_misconception} ($[0,1]^2$ is open in itself, not in $\mathbb{R}^2$), and it
returns as \cref{rem:boundedness_not_topological}. Compactness, by contrast, is
intrinsic like completeness: a subset $K$ is compact exactly when the metric space
$(K, d|_{K\times K})$ is, with no reference to the ambient space at all --- which is precisely
the reduction used in \cref{def:sequential_topological_compactness}.
\end{ainote}
\end{remark}

% Quelle: Diego Torres Tejeda/Notes - 06.03 - Diego Torres Tejeda.pdf, p. 2
\begin{exercise}[Completeness is not a topological property]
\label{ex:completeness_not_topological}
Corsin does not address whether completeness survives a change of metric; Diego Torres Tejeda
sets this exercise, which settles it. Let $X := (0,1]$, let $d_{\text{eucl}}(x,y) := |x-y|$ be
the standard metric on $X$, and define
\[d(x,y) := \left|\frac1x - \frac1y\right|, \qquad x,y \in X.\]
\begin{enumerate}[label=\textbf{(\alph*)}]
  \item Show that $d$ is a metric on $X$.
  \item Show that $(X,d_{\text{eucl}})$ is \textbf{not} complete.
  \item Show that $(X,d)$ \textbf{is} complete.
  \item Show that $d_{\text{eucl}}$ and $d$ induce the \textbf{same topology} on $X$.
\end{enumerate}
\textit{Diego's hint for} \textbf{(c)}: let $(x_n)$ be $d$-Cauchy. Then $(1/x_n)$ is Cauchy for
the standard metric, so it has a limit $\ell \in \mathbb{R}$; show $\ell \geq 1$ and that
$x_n \to 1/\ell$.
\end{exercise}

\begin{ainote}
Diego states the exercise on $X = (0,\infty)$, but there part \textbf{(c)} is false: on
$(0,\infty)$ the map $x\mapsto1/x$ is a bijection \emph{onto} $(0,\infty)$, so $(X,d)$ is
isometric to $\big((0,\infty),|\cdot|\big)$, which is not complete --- concretely $x_n := n$
satisfies $d(x_n,x_m) = |\tfrac1n-\tfrac1m| \to 0$ yet has no limit in $X$. His own hint asks to
show $\ell > 0$, and that is precisely the step that fails. Restricting to a bounded interval
repairs it: on $(0,1]$ the image is $[1,\infty)$, which \emph{is} closed in $\mathbb{R}$ and
hence complete. Typeset with $X := (0,1]$ above.
\end{ainote}

% Sonnet 5 (Medium)
Compactness will turn out to be a strictly stronger finiteness-type property than completeness:
every compact metric space is complete (\cref{ex:compact_finite_complete}), but the converse
fails --- $\mathbb{R}^n$ itself is complete, yet (being unbounded) not compact, as we will see
next (\cref{thm:sequential_compactness}, \cref{sec:heine_borel}).

% Source: Corsin Nick/Class Notes/Week 2.pdf, pp. 1--13
